\newpage
\section{Test Funzionali}

Un primo tipo di \textit{criteri di adeguatezza} di interesse sono i \textit{test funzionali}, che prendono questo nome in quanto il criterio di adeguatezza viene definito analizzando ``manualmente'' la specifica e trasformandola in un elenco di test obligations.

Tipicamente, tutti i criteri di adeguatezza di tipo funzionale (e quindi basati sulla specifica) cercano di riprodurre la specifica come una lista di comportamenti rispetto ai quali è possibile identificare: \begin{itemize}
    \item comportamenti normali del sistema;
    \item comportamenti eccezionali;
    \item casi di confine.
\end{itemize}

In generale:

\begin{figure}[H]
\caption{Tipi di approccio}
\begin{center}
\begin{tabular}{p{7cm} | p{7cm}}
    Approccio Sistematico & Approccio Random \\ \hline
    \begin{itemize}
        \item Cerca di fare delle scelte specifiche guidate da dei principi di analisi della specifica
        \item È legato alla persona che compie le scelte
    \end{itemize}
    & \begin{itemize}
        \item È più automatico
        \item Evita che la scelta dei casi di test sia basata sull'abilità specifica della persona che esegue l'analisi
        \item Sceglie una distribuzione per gli input e campiona automaticamente
    \end{itemize}
\end{tabular}
\end{center}
\end{figure}

Chiaramente, selezionare il tipo di approccio da usare di volta in volta è essenziale.

\begin{example}{}{Quando applicare un approccio sistematico?}
    Si supponga di dover testare un programma che risolve un'equazione di secondo grado, che sappiamo essere risolvibile con la seguente formula:
    $$x_{1,2} = \frac{- b \pm \sqrt{b^2 - 4ac}}{2a}$$
    L'input saranno i coefficienti $a,b,c$ ed il programma ritorna $x_{1,2}$.

    In questo caso, andare a testare esempi random non è l'approccio più corretto da fare, poiché non si vanno ad esaminare i casi più di interesse, come quando $\Delta = \sqrt{b^2 - 4ac} = 0$, oppure $2a = 0$.
\end{example}

\tikzset{
    boxtest/.style={draw, minimum width=1.5em, minimum height=1.5em, outer sep=0pt},
    boxtestfull/.style={boxtest, fill=red}
}

\begin{figure}[H]

\caption{Distribuzione delle failure}
\begin{center}

\begin{itemize}
    \item[] \begin{tikzpicture}\draw node[boxtestfull] {\;};\end{tikzpicture} Failure (test case utile)
    \item[] \begin{tikzpicture}\draw node[boxtest] {\;};\end{tikzpicture} No Failure 
\end{itemize}

\begin{tikzpicture}[node distance=1cm]
    \matrix (m) [matrix of nodes,nodes={boxtest},row sep=0.5em,column sep=0.5em]
    {
        \; & \; & & \; & \; & & \; & \; & & \; & \; & & |[boxtestfull]| \; & |[boxtestfull]| \\
        \; & \; & & \; & \; & & \; & \; & & \; & \; & & |[boxtestfull]| \; & \; \\
        \\
        \; & \; & & \; & \; & & \; & \; & & \; & \; & & \; & \; \\
        \; & \; & & \; & \; & & \; & \; & & \; & \; & & \; & \; \\
        \\
        \; & \; & & \; & \; & & \; & \; & & \; & \; & & \; & \; \\
        \; & \; & & \; & \; & & \; & \; & & \; & \; & & \; & \; \\
        \\
        \; & \; & & \; & \; & & |[boxtestfull]| \; & |[boxtestfull]| \; & & \; & \; & & \; & \; \\
        \; & \; & & \; & \; & & |[boxtestfull]| \; & |[boxtestfull]| \; & & \; & \; & & \; & \; \\
        \\
        \; & \; & & \; & \; & & \; & \; & & \; & \; & & \; & \; \\
        \; & \; & & \; & \; & & \; & \; & & \; & \; & & \; & \; \\
    };
\end{tikzpicture}
\end{center}

\end{figure}

Nella figura sopra, vediamo come le failure siano scarse in tutto il range di input, ma si addensino attorno a certi punti, che sono i nostro punti di interesse.

In generale, si vogliono eseguire poche prove, ma fondamentali. Ovviamente, è impossibile stabilire a priori, dato che gli input sono infiniti, dove si verificheranno esattamente le failure. Possiamo, però, dividere gli input possibili in insiemi, e per ciascuno di questi testare almeno un input, in maniera tale da poter coprire la maggior parte dei casi; alternativamente, si può ragionare a partire dal codice, supponendo che ogni statement corrisponda ad una classe di input. Questi, ed altri criteri per la scelta dei test, vengono definiti \textit{criteri di adeguatezza}.

\subsection{Criteri di adeguatezza}

Esistono diversi tipi di criteri di adeguatezza, ed ogni tipologia ha un'origine diversa.

\textbf{Criteri funzionali}: \\
derivano dalle specifiche del software

\textbf{Criteri strutturali}: \\
derivano dall'analisi del codice

\textbf{Criteri basati sui modelli}: \\
i modelli vengono estratti dalle specifiche (test su un protocollo di comunicazione)

\textbf{Criteri basati sui fault}: \\
si ipotizza che il programma possa contenere certe classi di fault (come per esempio il buffer overflow)

\subsubsection{Funzionali vs Strutturali}

In generale, nessuno dei due criteri prevale sull'altro, infatti spesso si adotta una combinazione dei due.

\begin{example}{}{Divisione tra numeri}
    Supponiamo di avere il seguente frammento di codice:

    \begin{lstlisting}
int fun (int param) {
    int result;
    result = param / 2;
    return result;
}
    \end{lstlisting}
    Ed analizziamolo sia dal punto di vista della specifica, che dal punto di vista del codice:

    \textbf{Specifica}: un programma che prende in input un intero \textit{n} e restituisce \textit{n / 2} se l'input è un numero pari, o solo \textit{n} se l'input è un numero dispari.

    \begin{figure}[H]
    \caption{Criteri a confronto}
    \begin{center}
    \begin{tabular}{p{5cm} | p{5cm}}
        Criterio basato sulla specifica & Criterio basato sul codice \\ \hline
        Sono necessari almeno due test: \begin{itemize}
            \item Almeno un test per input pari
            \item Almeno un test per input dispari
        \end{itemize}
        & È sufficiente un unico test per eseguire tutti gli statement
    \end{tabular}
    \end{center}
    \end{figure}

    Nel caso del \textit{criterio basato sulla specifica}, si rivelerà un malfunzionamento durante il test con input dispari.

    Con il \textit{criterio basato sul codice}, non avrò la certezza di individuare il bug, poiché potrebbe essere testato un singolo numero pari.

    Quindi è chiaro, in questo caso, che la scelta migliore sia di adottare un \textit{criterio basato sul codice}; si può inoltre notare come in questo esempio venga evidenziata la mancanza di implementazione di una parte della specifica (errore di \textit{missing logic}).

    In generale, il \textit{criterio basato sulla specifica} ci aiutare a trovare facilmente errori di \textit{missing logic}.
\end{example}

\subsection{Combinational Testing}

Una tecnica di testing funzionale è il \textit{combinational testing}, che è identificata da due step: \begin{itemize}
    \item passo di decomposizione \\ il problema viene partizionato in maniera da poter essere più facilmente gestibile. La specifica viene trasformata in una lista di elementi (\textit{test obligations}) che la rappresentano;
    \item passo di aggregazione \\ i dati estratti precedentemente vengono aggregati (o \textit{combinati}, termine che da il nome alla tecnica) in maniera tale da ottenere dei casi di test intelligenti oltre che completi.
\end{itemize}

Per ragionare su come ottenere il risultato suggerito da questi due step, useremo tre tecniche: \begin{itemize}
    \item Category-partition testing \\ si cerca di risolvere il problema della scomposizione della specifica nei suoi valori interessanti
    \item Pairwise testing \\ ricombina tutti i possibili input per cercare di risolvere il problema dell'esplosione combinatoria
    \item Catalog-based testing \\ approccio più sistematico che può essere integrato con l'analisi della specifica, si cerca che tipo di valori si devono andare a cercare
\end{itemize}

\begin{figure}[H]
\caption{Tipi di test ed approccio}
\begin{center}
\begin{tabular}{p{5cm} | p{5cm}}
    Passo di decomposizione & Passo di aggregazione \\ \hline
    \begin{itemize}
        \item Category-partition testing
        \item Catalog-based testing
    \end{itemize}
    & \begin{itemize}
        \item Pariwise testing
    \end{itemize}
\end{tabular}
\end{center}
\end{figure}

\subsubsection{Pairwise Testing}

Lo scopo di questo step non è solamente considerare la lista pura dei comportamenti estratta dalla specifica, ma piuttosto le combinazioni di questi comportamenti. Questo è quindi il \textit{pairwise testing}.

\begin{example}{}{Display Control}
    Si immagini di avere cinque parametri (1 colonna = 1 parametro del sistema). La specifica non ci è data a sapere, ma abbiamo già estratto tutti i valori significativi. Questi parametri sono stati rivelati in fase di test. 
\end{example}